\section{Feature Importance and Interpretability}
\label{sec:interpretability}

\subsection{Random Forest Importance Analysis}
\label{sec:interp_rf}
Understanding which descriptor families drive predictions supports clinical
interpretability and validates Phase~1 separability rankings. We aggregate
\texttt{feature\_importances\_} from the best Random Forest configuration on
$\mathrm{Full(opt)}$ by feature group: statistics ($11$ dims), GLCM ($6$),
LBP ($14$), DWT ($16$), and HOG ($\sim$490 after cleaning).

\subsection{Feature Group Contribution}
\label{sec:interp_groups}
Figure~\ref{fig:rf_importance} shows RF cumulative group importance. HOG
dominates raw importance owing to its dimensionality; per-dimension importance
ranks GLCM and statistical descriptors higher, consistent with their compact
yet discriminative encoding of intensity and texture structure.

\begin{figure}[t]
\centering
\includegraphics[width=\linewidth]{rf_branch_importance}
\caption{Random forest feature-group importance on $\mathrm{Full(opt)}$.}
\label{fig:rf_importance}
\end{figure}

\subsection{Linking Phase 1 and Phase 2}
\label{sec:interp_link}
High Phase~1 separability scores for LBP, DWT, and HOG align with strong
Phase~2 performance when those blocks are used alone or in fusion ($D$ mean
F1 $=0.865$; $\mathrm{Full(opt)}$ F1 $=0.920$ with SVM). This pattern
supports the classifier-free selection strategy: separability metrics identify
informative descriptors before expensive classifier sweeps.
